% !TeX encoding = UTF-8
% !TeX program = pdflatex
% ---------------------------------------------------------------------------
% IEEE TIP (Transactions on Image Processing) Journal Template
%
% Usage:
%   1. Requires IEEEtran.cls and IEEEtran.bst (copy from templates/ieee-latex/)
%   2. Compile: pdflatex → bibtex → pdflatex → pdflatex
%
% Journal: IEEE Transactions on Image Processing
% Publisher: IEEE Signal Processing Society
% ISSN: 1057-7149
% Impact Factor: ~10.6 (2024 JCR)
% Scope: Image/video processing, computer vision, image analysis,
%        image coding, restoration, enhancement, and related theory
% ---------------------------------------------------------------------------

\documentclass[journal,onecolumn,draftcls]{IEEEtran}
\IEEEoverridecommandlockouts

% --- Final version ---
% \documentclass[journal]{IEEEtran}
% \IEEEoverridecommandlockouts

\usepackage{cite}
\usepackage{amsmath,amssymb,amsfonts}
\usepackage{algorithmic}
\usepackage{algorithm}
\usepackage{graphicx}
\usepackage{textcomp}
\usepackage{xcolor}
\usepackage{booktabs}
\usepackage{multirow}
\usepackage{hyperref}
\usepackage{subcaption}
\usepackage{adjustbox}

\hypersetup{
  colorlinks=true,
  linkcolor=blue,
  citecolor=blue,
  urlcolor=blue,
}

\def\BibTeX{{\rm B\kern-.05em{\sc i\kern-.025em b}\kern-.08em
    T\kern-.1667em\lower.7ex\hbox{E}\kern-.125emX}}

% --- TIP Journal Metadata ---
% \markboth{IEEE Transactions on Image Processing, Vol. XX, No. X, Month 2026}%
% {Author \MakeLowercase{\textit{et al.}}: Short Title Here}

\begin{document}

% ===========================================================================
% TITLE
% ===========================================================================
\title{Your Paper Title: [Method Name] for [Image Processing / Computer Vision Task]}

\author{
  \IEEEauthorblockN{Author One\IEEEauthorrefmark{1},
                    Author Two\IEEEauthorrefmark{2},
                    Author Three\IEEEauthorrefmark{1}}
  \IEEEauthorblockA{
    \IEEEauthorrefmark{1}Department of XXX, University of YYY, City, Country \\
    \IEEEauthorrefmark{2}Institute of ZZZ, Academy of Sciences, City, Country \\
    E-mail: \{email1, email2\}@university.edu
  }
}

\maketitle

% ===========================================================================
% ABSTRACT
% ===========================================================================
\begin{abstract}

% TIP abstracts: 150-250 words, single paragraph.
% Must clearly convey the image processing contribution.
%
% TIP emphasizes: (1) image/video processing novelty,
% (2) rigorous experimental validation on standard benchmarks,
% (3) theoretical justification where applicable.

Your abstract here. Structure:
\begin{itemize}
  \item \textbf{Problem:} What image/video processing task and why it matters.
  \item \textbf{Gap:} Why existing image processing methods are insufficient.
  \item \textbf{Method:} First... Second... Finally... (core IP/CV contribution).
  \item \textbf{Results:} Quantitative evidence on standard benchmarks.
  \item \textbf{Significance:} Broader impact for image processing community.
\end{itemize}

\end{abstract}

% ===========================================================================
% KEYWORDS
% ===========================================================================
\begin{IEEEkeywords}
  Image processing, [your task, e.g., image restoration, super-resolution,
  denoising, segmentation], [your technique, e.g., deep learning, attention
  mechanism, diffusion models], [secondary technique],
  [application domain if applicable].
\end{IEEEkeywords}

% ===========================================================================
% 1. INTRODUCTION
% ===========================================================================
\section{Introduction}
\label{sec:introduction}

% TIP Introduction (5 paragraphs, ~2-3 pages in journal submission):

% Para 1: Image processing context and importance.
%   - Why this image/video processing task matters.
%   - Real-world applications (medical imaging, photography, autonomous
%     driving, surveillance, remote sensing, entertainment, etc.).
%   - Connect to the broader signal/image processing community.

% Para 2: Technical progress.
%   - Evolution: classical signal processing → hand-crafted features →
%     deep learning → foundation models / diffusion / transformers.
%   - 2-3 major method families with representative works.
%   - Cite TIP, TIP, TIP papers to show you know the venue.

% Para 3: Remaining gap — must have a "because" clause.
%   - Image-processing-specific gap examples:
%     "Existing denoising methods assume AWGN, but real-world sensor noise
%      follows signal-dependent distributions..."
%     "Current SR methods produce plausible but unfaithful details because..."
%     "Segmentation models fail on fine structures because..."
%   - Quantify the gap if possible (cite numbers from prior work).

% Para 4: Proposed method teaser (3-4 sentences).
%   - Core image processing insight.
%   - WHY it addresses the gap at the signal/image level.

% Para 5: Main contributions (4-5 bullets).

The main contributions of this paper are:
\begin{enumerate}
  \item We identify [specific limitation] in existing [method family] for
    [image processing task], validated through [analysis / preliminary
    experiments] (Section~\ref{sec:motivation}).
  \item We propose [Method Name], which [core image processing mechanism]
    to address [specific challenge] (Section~\ref{sec:method}).
  \item We provide [theoretical analysis / mathematical justification]
    showing that [property] holds, and we analyze [complexity / convergence /
    optimality] (Section~\ref{sec:theory}).
  \item Extensive experiments on [N] benchmark datasets demonstrate that
    [Method] achieves [key quantitative result], outperforming [M] SOTA
    methods, with [specific advantage, e.g., "3× fewer parameters" or
    "2.1 dB PSNR gain"] (Section~\ref{sec:experiments}).
  \item Code and pretrained models are available at [URL].
\end{enumerate}

% ===========================================================================
% 2. RELATED WORK
% ===========================================================================
\section{Related Work}
\label{sec:related}

% TIP Related Work: comprehensive, 3-4 pages in journal mode.
% Organize by method family AND by task.

\subsection{Image [Task]: Classical Approaches}
\label{sec:related-classical}

% Classical signal/image processing approaches:
% - Filtering-based methods (bilateral, guided, BM3D, NLM)
% - Variational / PDE-based methods (TV, anisotropic diffusion)
% - Sparse representation / dictionary learning
% - Wavelet / frequency-domain methods
% Show you understand the fundamentals of image processing, not just DL.

\subsection{Deep Learning for Image [Task]}
\label{sec:related-dl}

% Deep learning evolution: CNN → ResNet/UNet → Transformer/ViT →
% Diffusion → Mamba/SSM.
% Organize by architecture type, not paper-by-paper.

\subsection{[Method Paradigm A]}
\label{sec:related-paradigmA}

% e.g., GAN-based methods, diffusion-based methods, NeRF-based methods

\subsection{[Method Paradigm B]}
\label{sec:related-paradigmB}

% e.g., Self-supervised, multi-task, multi-modal approaches

\subsection{Summary and Positioning}
\label{sec:related-summary}

% Explicitly state what ALL prior work misses and how your method differs.
% TIP reviewers want a clear contribution statement in Related Work.

% ===========================================================================
% 3. PRELIMINARIES
% ===========================================================================
\section{Preliminaries}
\label{sec:preliminaries}

\subsection{Problem Formulation}
\label{sec:prelim-problem}

% Formal mathematical definition of the image processing task.
% Example for image restoration:
%   $\mathbf{y} = \mathbf{H}\mathbf{x} + \mathbf{n}$
%   where $\mathbf{y}$ is the degraded observation,
%   $\mathbf{H}$ is the degradation operator,
%   $\mathbf{x}$ is the clean image,
%   $\mathbf{n}$ is noise.
% Goal: recover $\hat{\mathbf{x}} \approx \mathbf{x}$ from $\mathbf{y}$.

\subsection{Background Concepts}
\label{sec:prelim-bg}

% Define key concepts needed to understand your method.
% Keep this brief -- TIP readers are image processing experts.

% ===========================================================================
% 4. PROPOSED METHOD
% ===========================================================================
\section{Proposed Method}
\label{sec:method}

% TIP method sections: very detailed (3-5 pages).
% A reader should be able to re-implement from this description.
% Every equation must be explained; every architecture choice justified.

\subsection{Overview}
\label{sec:method-overview}

% Architecture overview diagram (Fig. 1).
% Describe the complete pipeline from input image to output.

\begin{figure}[t]
  \centering
  \fbox{\parbox{0.85\linewidth}{\centering
  \vspace{3cm}
  {\large [Fig. 1: Architecture overview of the proposed method.]}\\
  \vspace{0.3cm}
  {\footnotesize The pipeline: degraded image → [preprocessing] →
   [feature extraction] → [core module] → [reconstruction] → restored image.}
  \vspace{3cm}}}
  \caption{Overview of the proposed [Method]. [One sentence describing data flow].}
  \label{fig:architecture}
\end{figure}

\subsection{[Core Module A]}
\label{sec:method-moduleA}

% For each module: Input → Process → Output → Rationale.
% Image processing specifics:
% - Specify tensor shapes: (B, C, H, W) for images, (B, T, C, H, W) for video
% - Define kernel sizes, stride, padding for convolution layers
% - Define window size, patch size for transformer-based methods
% - Define noise level, degradation model for restoration methods

\subsection{[Core Module B]}
\label{sec:method-moduleB}

% Second innovation.

\subsection{Loss Function Design}
\label{sec:method-loss}

% TIP papers often have carefully designed loss functions.
% Common image processing losses:
% - Reconstruction: L1 (MAE), L2 (MSE), Charbonnier
% - Perceptual: VGG perceptual loss, LPIPS
% - Adversarial: GAN loss (LSGAN, hinge, WGAN-GP)
% - Structural: SSIM loss, MS-SSIM loss
% - Frequency: FFT loss, wavelet loss
% - Edge-aware: gradient loss, Laplacian loss

% Explain WHY each loss term is included and what image property it preserves.

\begin{equation}
  \mathcal{L}_{\text{total}} =
    \lambda_1 \mathcal{L}_{\text{reconstruction}} +
    \lambda_2 \mathcal{L}_{\text{perceptual}} +
    \lambda_3 \mathcal{L}_{\text{[auxiliary]}}
  \label{eq:total_loss}
\end{equation}

% Each λ chosen via [validation / grid search / adaptive weighting].

\subsection{Implementation Details}
\label{sec:method-implementation}

% Comprehensive for TIP reproducibility:
% - Framework and version (PyTorch X.X, CUDA X.X)
% - Optimizer, learning rate, scheduler
% - Batch size, patch size, number of epochs/iterations
% - Data augmentation (geometric, photometric, degradation-specific)
% - Hardware and training time
% - Pretrained weights (if used)

% ===========================================================================
% 5. EXPERIMENTS
% ===========================================================================
\section{Experiments}
\label{sec:experiments}

% TIP experiments: very comprehensive, 4-6 pages.
% Must include: datasets, metrics, main comparison, ablation, efficiency,
% qualitative visual comparison, failure case analysis.

% ---------------------------------------------------------------------------
\subsection{Datasets and Evaluation Protocols}
\label{sec:exp-datasets}

% For each dataset: source, size, resolution, characteristics, split.
% Standard image processing benchmarks:

\begin{table}[t]
  \centering
  \caption{Summary of datasets used in experiments.}
  \label{tab:datasets}
  \begin{tabular}{lccrl}
    \toprule
    \textbf{Dataset} & \textbf{Images} & \textbf{Resolution} &
    \textbf{Task} & \textbf{Split (Train/Val/Test)} \\
    \midrule
    DIV2K & 1,000 & 2K & SR & 800/100/100 \\
    Flickr2K & 2,650 & 2K & SR & 2,450/0/200 \\
    BSD68 & 68 & 481×321/321×481 & Denoising & --/--/68 \\
    Set5/Set14 & 5/14 & various & SR/Denoising & --/--/5,14 \\
    Urban100 & 100 & various & SR & --/--/100 \\
    \bottomrule
  \end{tabular}
\end{table}

% ---------------------------------------------------------------------------
\subsection{Evaluation Metrics}
\label{sec:exp-metrics}

% TIP standard metrics:
% - Full-reference: PSNR, SSIM, MS-SSIM, LPIPS, DISTS, NLPD, FSIM, VIF
% - No-reference: NIQE, BRISQUE, PIQE, MANIQA, CLIPIQA
% - Task-specific: mIoU (segmentation), AP (detection), FID (generation)
% - Efficiency: Params (M), FLOPs (G), runtime (ms), GPU memory (GB)
%
% State metric direction: higher is better for PSNR/SSIM/PSNR;
% lower is better for LPIPS/NIQE/runtime.

% ---------------------------------------------------------------------------
\subsection{Comparison with State-of-the-Art}
\label{sec:exp-main}

% Table: method, year, venue, PSNR, SSIM, LPIPS, Params, FLOPs, Runtime.
% Bold = best, underline = second-best.
% Multiple scale factors (×2, ×3, ×4) in separate tables or sub-tables.
% If reporting on multiple noise levels, organize by σ.

\begin{table}[t]
  \centering
  \caption{Quantitative comparison on [Dataset] for [Task].
           Best results in \textbf{bold}, second-best \underline{underlined.}}
  \label{tab:main}
  \begin{tabular}{lcccccc}
    \toprule
    \textbf{Method} & \textbf{PSNR↑} & \textbf{SSIM↑} & \textbf{LPIPS↓} &
    \textbf{Params (M)} & \textbf{FLOPs (G)} & \textbf{Runtime (ms)} \\
    \midrule
    Method A \cite{ref1} & 32.10 & 0.895 & 0.112 & 15.2 & 88.3 & 45 \\
    Method B \cite{ref2} & 32.85 & 0.902 & 0.095 & 32.0 & 156.7 & 78 \\
    \textbf{Ours} & \textbf{33.21} & \textbf{0.910} &
    \textbf{0.078} & 18.5 & 92.1 & 38 \\
    \bottomrule
  \end{tabular}
\end{table}

% ---------------------------------------------------------------------------
\subsection{Ablation Study}
\label{sec:exp-ablation}

% Comprehensive ablation of ALL components.
% Waterfall-style: Full model → remove A → remove B → remove both → baseline.
% Show PSNR drop for each removal.

% ---------------------------------------------------------------------------
\subsection{Model Efficiency Analysis}
\label{sec:exp-efficiency}

% TIP highly values efficiency -- image processing needs to be practical.
% - Parameter count comparison (bar chart)
% - FLOPs vs PSNR scatter (Pareto frontier)
% - Runtime on different resolutions
% - GPU memory consumption
% - Throughput (FPS) for video tasks

% ---------------------------------------------------------------------------
\subsection{Qualitative Results}
\label{sec:exp-qualitative}

% Visual comparison with SOTA methods.
% Show zoomed-in patches to highlight fine details.
% Use consistent color mapping across methods.
% Show at least 2-3 diverse examples per dataset.
% Include both easy and challenging cases.

% ---------------------------------------------------------------------------
\subsection{Failure Case Analysis}
\label{sec:exp-failure}

% TIP expects honesty about limitations.
% Show cases where the method fails:
% - What types of images/scenes cause degradation?
% - What image structures are poorly reconstructed?
% - Is there a systematic bias (e.g., over-smoothing, hallucination)?
% Discuss WHY these failures occur.

% ---------------------------------------------------------------------------
\subsection{Generalization Analysis}
\label{sec:exp-generalization}

% TIP values generalization:
% - Cross-dataset: train on A, test on B
% - Cross-degradation: train on one noise level, test on others
% - Cross-resolution: train on 2K, test on 4K images
% - Real-world: test on real (not synthetic) degraded images

% ===========================================================================
% 6. DISCUSSION
% ===========================================================================
\section{Discussion}
\label{sec:discussion}

\subsection{Interpretation of Results}
\label{sec:discuss-interpretation}

% Explain WHY your method outperforms baselines.
% Connect experimental observations to design choices.

\subsection{Limitations}
\label{sec:discuss-limitations}

% Honest assessment of limitations.
% Image-processing-specific:
% - Resolution dependence
% - Content dependence (textures, faces, text)
% - Sensor/noise-type dependence
% - Computational constraints for real-time applications

\subsection{Future Work}
\label{sec:discuss-future}

% Concrete image processing research directions.

% ===========================================================================
% 7. CONCLUSION
% ===========================================================================
\section{Conclusion}
\label{sec:conclusion}

% ---------------------------------------------------------------------------
% 8. ACKNOWLEDGMENTS
% ---------------------------------------------------------------------------
\section*{Acknowledgments}

% ===========================================================================
% REFERENCES
% ===========================================================================
{
  \bibliographystyle{IEEEtran}
  \bibliography{references}
}

\end{document}
